\documentclass[12pt]{article}
\usepackage[utf8]{inputenc}
\usepackage[T1]{fontenc}
\usepackage{hyperref}
\title{Tajio Language Resources for K-12 Curriculum Development}
\author{Compiled by Claude Sonnet 5 (claude-sonnet-5) for Andr\'as Kornai}
\date{}
\begin{document}
\maketitle
\section{Introduction}
Tajio (also Kasimbar, Tadjio, Ajio; ISO 639-3 tdj; Glottocode taji1246) is a Celebic Austronesian language of Central Sulawesi, Indonesia. It is spoken mainly in Donggala Regency. Estimates range from 8,800 to 12,000 speakers. Most speakers use Indonesian, and many also use Kaili.
\section{Primary Educational Resources}
No school materials were found in search results.
\section{Linguistic and Reference Materials}
\begin{itemize}
\item Mayani, Luh Anik: PhD grammar of Tajio (Universit\"at zu K\"oln, 2013): \url{https://kups.ub.uni-koeln.de/6963/}
\item Tajio Webonary (dictionary site): \url{https://www.webonary.org/tajio/overview/introduction/}
\item Omniglot page on Tajio: \url{https://omniglot.com/writing/tajio.htm}
\item Wikipedia overview: \url{https://en.wikipedia.org/wiki/Tajio_language}
\end{itemize}
\section{Religious and Community Resources}
Joshua Project (reports some Bible portions): \url{https://joshuaproject.net/people_groups/15167}
\section{Digital and Audio Resources}
The Webonary site above is the main digital resource. Endangered Languages Project entry: \url{https://endangeredlanguages.com/lang/2786}
\section{Curriculum Development Context}
Tajio has a Latin-script writing system, a full grammar and an online dictionary. These are useful foundations. No existing school programme was found.
\end{document}